% Compile twice with local XeLaTeX:
% xelatex -interaction=nonstopmode -halt-on-error proposal_defense.tex
% Content source: 许哲圣_本科毕业论文开题报告(1111).docx
% Five slides. Speaker script: proposal_speech.tex / proposal_speech.pdf
\documentclass[aspectratio=169,12pt]{beamer}
\usepackage[UTF8,fontset=fandol]{ctex}
\usepackage{amsmath,amssymb}

\definecolor{Navy}{HTML}{17385B}
\definecolor{TextGray}{HTML}{374151}
\setbeamercolor{normal text}{fg=TextGray,bg=white}
\setbeamercolor{structure}{fg=Navy}
\setbeamercolor{frametitle}{fg=Navy,bg=white}
\setbeamercolor{title}{fg=Navy}
\setbeamercolor{itemize item}{fg=Navy}
\setbeamerfont{frametitle}{size=\Large,series=\bfseries}
\setbeamerfont{title}{size=\LARGE,series=\bfseries}
\setbeamertemplate{navigation symbols}{}
\setbeamertemplate{itemize item}{\small\textbullet}
\setbeamertemplate{footline}{%
  \hfill\usebeamercolor[fg]{structure}\scriptsize
  \insertframenumber\,/\,\inserttotalframenumber\hspace{6mm}\vspace{4mm}}
\setbeamersize{text margin left=10mm,text margin right=10mm}
\setlength{\parskip}{5pt}
\hypersetup{pdftitle={基于宽松风险平价的多资产ETF配置研究 开题答辩},pdfauthor={许哲圣}}

\begin{document}

% Slide 1: 20 seconds
\begin{frame}
  \vspace{7mm}
  {\color{Navy}\small 本科毕业论文开题答辩\par}
  \vspace{6mm}
  {\color{Navy}\LARGE\bfseries 基于宽松风险平价的\par 多资产 ETF 配置研究\par}
  \vspace{8mm}
  {\normalsize 许哲圣\qquad 指导教师：王磊\par}
  {\small 特拉华数据科学学院\quad 金融数学专业\par}
  \vspace{2mm}
  {\small 2026 年 11 月\par}
\end{frame}

% Slide 2: 65 seconds
\begin{frame}{研究背景与研究问题}
  \begin{itemize}
    \setlength{\itemsep}{10pt}
    \item 股票、债券等资产波动不同，资金等权不等于风险均衡。
    \item 风险预算提供配置参考，但仍需权衡总风险与预期收益。
    \item 允许组合在参考配置附近调整，检验调整是否值得。
  \end{itemize}
  \vspace{6mm}
  {\color{Navy}\bfseries 核心问题\par}
  在风险预算参考基础上加入方差和预测收益短缺惩罚，\par
  能否在扣除交易成本后形成合理的收益与风险取舍？
\end{frame}

% Slide 3: 120 seconds
\begin{frame}{两阶段模型与核心公式}
  \textbf{第一阶段：}求解等风险预算组合，得到参考权重 $q$。\par
  \textbf{第二阶段：}在参考组合附近优化最终权重 $w$。
  \vspace{1mm}
  \[
    \min_w\quad
    \underbrace{\lVert w-q\rVert_2^2}_{\text{偏离参考}}
    +\underbrace{\lambda_v\frac{w^{\mathsf T}\Sigma w}{v}}_{\text{组合方差}}
    +\underbrace{\lambda_r\!\left[\max\!\left(0,\frac{R-\mu^{\mathsf T}w}{s}\right)\right]^2}_{\text{预测收益短缺}}
  \]
  \[
    \mathbf{1}^{\mathsf T}w=1,\qquad w_i\geq 0,\qquad R=\mu^{\mathsf T}q.
  \]
  {\small
    $\mu$ 为预测年化收益，$\Sigma$ 为年化协方差矩阵。\par
    $v,s>0$ 为求解前固定的尺度，$\lambda_v,\lambda_r\geq0$ 为惩罚系数。\par
  }
  \vspace{2mm}
  {\color{Navy}“宽松”指允许最终配置偏离严格的等风险预算关系。}
\end{frame}

% Slide 4: 75 seconds
\begin{frame}{数据与验证方法}
  \begin{itemize}
    \setlength{\itemsep}{9pt}
    \item \textbf{研究对象：}30 只境内上市 ETF，覆盖股票、债券、\newline
      货币市场及商品等类别。
    \item \textbf{回测安排：}使用 2018 年至 2026 年的历史数据，\newline
      周度调仓，计入单边 3 个基点交易成本。
    \item \textbf{信息使用：}仅使用历史数据估计输入、选择年度参数，\newline
      参数选定后在下一年度保持固定。
    \item \textbf{效果检验：}对比风险预算参考及简单配置，去掉单项\newline
      惩罚作比较，结合收益、波动、回撤与持仓解释结果。
  \end{itemize}
\end{frame}

% Slide 5: 10 seconds
\begin{frame}
  \centering
  \vspace{10mm}
  {\color{Navy}\LARGE\bfseries 感谢各位老师\par}
  \vspace{6mm}
  {\large 请老师批评指正\par}
\end{frame}

\end{document}
